% Author: Writer agent · 2026-11-14
% Task W2: Chapter 8 (epilogue). Sources: research-notes.md, Ch. 8 (recap keys reused).

\chapter{נועה לוחצת Play}\label{ch:8}

\begin{scene}
07:30, שוב. השעון המעורר, מברשת השיניים, מיקס הבוקר. הכול כמו בפרק~\ref{ch:1}, חוץ מדבר אחד: הפעם נועה יודעת מה קורה מאחורי המסך. הדילוג באוטובוס הוא רמז, לא פסק דין. יוני הוא וקטור שכן. השיר החדש הגיע אולי דרך השמע שלו. ואם יתנגן הבוקר משהו מוזר לגמרי, אולי זו פשוט חקירה.

היא מחייכת אל הטלפון, מוחקת את קצף משחת השיניים מהסנטר, ולוחצת \en{Play}.
\end{scene}

\section{למה האותות שלך חשובים}

לפי \en{Spotify}, פרופיל הטעם של המשתמש הוא הקלט החשוב ביותר ליצירת חוויית השימוש הטובה ביותר \cite{spotify2026understanding}. במילים אחרות, מה שנועה עושה משפיע על מה שהיא תשמע מחר. הנה מה שלקחנו מהדרך:

\begin{inshort}
\begin{enumerate}
  \item \textbf{פרק 2:} משוב משתמע מבטא העדפה בדרגות ביטחון שונות \cite{hu2008collaborative}, ולכן דילוג יכול לומר „לא אהבתי", אבל גם „לא עכשיו".
  \item \textbf{פרק 3:} סינון שיתופי נשען על הרעיון שמי שהסכימו בעבר כנראה יסכימו גם בעתיד \cite{resnick1994grouplens}.
  \item \textbf{פרק 4:} המלצה מבוססת תוכן עוזרת לשירים חדשים מול בעיית ההתחלה הקרה, כי היא לומדת מהשיר עצמו \cite{vandenoord2013deep}.
  \item \textbf{פרק 5:} המערכת צריכה לאזן בין חקירה לניצול, בין הבטוח לבין החדש \cite{mcinerney2018explore}.
  \item \textbf{פרק 6:} הדפוס הכללי הוא יצירת מועמדים, דירוג, ואז סינון וגיוון \cite{google2026overview,covington2016deep}.
  \item \textbf{פרק 7:} האזנה אלגוריתמית קשורה לגיוון נמוך יותר, וגיוון גבוה קשור דווקא להישארות בשירות \cite{anderson2020algorithmic}.
\end{enumerate}
\end{inshort}

ואם המיקס של מחר יהיה קצת פחות מושלם? אולי זו לולאת המשוב, ואולי פשוט חקירה. ועכשיו נועה יודעת מה לשאול.
